%%%%%%%%%%%%%%%%%%%%%%%%%%%%%%%%%%%%%%%%%%%%%%%%%%%%%%%%%%%%%%%%%%%%%%%%%%%%%%%%
% Preámbulo                                                                    %
%%%%%%%%%%%%%%%%%%%%%%%%%%%%%%%%%%%%%%%%%%%%%%%%%%%%%%%%%%%%%%%%%%%%%%%%%%%%%%%%

\documentclass[11pt,a4paper,titlepage,oneside]{report}

%%% RELACIÓN DE VARIABLES A PERSONALIZAR %%%
\def\lingua{gal}
%\def\lingua{esp} % descomenta esta liña se redactarás a memoria en español
%\def\lingua{eng} % descomenta esta liña se redactarás a memoria en inglés
\def\nome{Francisco Lareo García}                             % substitúe aquí o teu nome
\def\nomedirectorA{Outro Nome Completo}              % substitúe aquí o nome de quen dirixe
%\def\nomedirectorB{Outro Nome Completo}             % duplica esta liña máis veces se o precisas, cambiando
                                                     % a letra final (A, B, C, D...): úsanse na portada.tex
\def\titulo{Desarrollo de una Aplicación Multiplataforma Escalable y Plataforma Analítica Complementaria
para la Recopilación y Análisis de Datos de Uso} % substitúe aquí o título do teu TFG
%\def\titulacion{gced}                               % descomenta esta liña e comenta a seguinte se es estudante do GCED
\def\titulacion{gei}
%\def\mencion{COMPUTACIÓN}                           % descomenta a mención que che corresponda se es estudante do GEI
\def\mencion{ENXEÑARÍA DO SOFTWARE}
%\def\mencion{ENXEÑARÍA DE COMPUTADORES}
%\def\mencion{SISTEMAS DE INFORMACIÓN}
%\def\mencion{TECNOLOXÍAS DA INFORMACIÓN}

%\def\renomearcadros{si} % descomenta esta liña se redactas a memoria en español e prefires que
                         % os "cuadros" e o "índice de cuadros" se renomeen
                         % a "tablas" e "índice de tablas" respectivamente

\usepackage{estilo_tfg}

% Lista de paquetes potencialmente interesantes (uso baixo demanda)

% \usepackage{alltt}       % proporciona o entorno alltt, semellante a verbatim pero que respecta comandos
% \usepackage{enumitem}    % permite personalizar os entornos de lista
% \usepackage{eurofont}    % proporciona o comando \euro
% \usepackage{float}       % permite máis opcións para controlar obxectos flotantes (táboas, figuras)
% \usepackage{hhline}      % permite personalizar as liñas horizontais en arrays e táboas
  \usepackage{longtable}   % permite construir táboas que ocupan máis dunha páxina
% \usepackage{lscape}      % permite colocar partes do documento en orientación apaisada
% \usepackage{moreverb}    % permite personalizar o entorno verbatim
  \usepackage{multirow}    % permite crear celdas que ocupan varias filas da mesma táboa
% \usepackage{pdfpages}    % permite insertar ficheiros en PDF no documento
% \usepackage{rotating}    % permite diferentes tipos de rotacións para figuras e táboas
% \usepackage{subcaption}  % permite a inclusión de varias subfiguras nunha figura
% \usepackage{tabu}        % permite táboas flexibles
% \usepackage{tabularx}    % permite táboas con columnas de anchura determinada

%%%%%%%%%%%%%%%%%%%%%%%%%%%%%%%%%%%%%%%%%%%%%%%%%%%%%%%%%%%%%%%%%%%%%%%%%%%%%%%%
% Corpo                                                                        %
%%%%%%%%%%%%%%%%%%%%%%%%%%%%%%%%%%%%%%%%%%%%%%%%%%%%%%%%%%%%%%%%%%%%%%%%%%%%%%%%

\begin{document}

 %%%%%%%%%%%%%%%%%%%%%%%%%%%%%%%%%%%%%%%%
 % Preliminares do documento            %
 %%%%%%%%%%%%%%%%%%%%%%%%%%%%%%%%%%%%%%%%

 \begin{titlepage}
  
  \hspace*{128pt}
  \textcolor{udcpink}{{\fontencoding{T1}\fontfamily{phv}\selectfont Facultade de Informática}}\\[-32pt]

  \begin{center}
    \includegraphics[scale=0.3]{imaxes/udc}\\[25pt]

    {\large TRABALLO FIN DE GRAO \\
            \nometitulacion \\
            \nomemencion } \\[10pt]

    \carimbo \\[25pt]

    \begin{huge}
      \begin{spacing}{1.3}
        \bfseries \titulo
      \end{spacing}
    \end{huge}
  \end{center}
  
  \vfill
  
  \begin{flushright}
    {\large
    \begin{tabular}{ll}
      {\bf Estudante:} & \nome \\
      {\bf Dirección:} & \nomedirectorA \\
%                      & \nomedirectorB \\ % duplica esta liña máis veces se o precisas, cambiando
                                           % a letra final (A, B, C, D...); define eses nomes no memoria_tfg.tex
    \end{tabular}}
  \end{flushright}
  \rightline{A Coruña, \datasimple.}
\end{titlepage}

 \dedicatoria{Dedicatoria} % escribe neste comando o teu texto de dedicatoria
 \paxinaenbranco
 \begin{agradecementos}
 \blindtext                % substitúe este comando polo teu texto de agradecementos
 \end{agradecementos}
 %%%%%%%%%%%%%%%%%%%%%%%%%%%%%%%%%%%%%%%%%%%%%%%%%%%%%%%%%%%%%%%%%%%%%%%%%%%%%%%%

\pagestyle{empty}
\begin{abstract}
  Este trabajo presenta el desarrollo e implementación de una plataforma OTT (Over-The-Top) multiplataforma, 
  diseñada para funcionar en una amplia gama de dispositivos, desde televisores inteligentes hasta ordenadores y 
  dispositivos móviles. La aplicación utiliza una arquitectura basada en microservicios, lo que garantiza escalabilidad, 
  flexibilidad y fácil mantenimiento. El objetivo principal del proyecto fue crear una plataforma robusta y adaptable, 
  capaz de integrarse con diversos sistemas operativos y satisfacer las necesidades de múltiples clientes.

A lo largo del proyecto, se han abordado importantes desafíos técnicos, como la adaptación de la aplicación a diferentes 
capacidades de procesamiento y la optimización del rendimiento en distintos entornos, tales como WebOS, Tizen y Android TV. 
Para ello, se han realizado pruebas exhaustivas, tanto a nivel unitario como de integración y sistema, que han permitido 
validar la funcionalidad y el rendimiento de la aplicación en condiciones reales de uso.

Además de la plataforma OTT, el proyecto incluye una aplicación de análisis de datos, desarrollada con la API de Matomo, 
que permite realizar un seguimiento detallado del comportamiento de los usuarios. Esta herramienta facilita la toma de 
decisiones basada en datos, permitiendo a los clientes obtener informes personalizados sobre el rendimiento de sus contenidos.

Los resultados obtenidos durante el desarrollo de este proyecto demuestran la capacidad de la plataforma para adaptarse a 
las demandas del mercado y proporcionar una experiencia de usuario óptima. Finalmente, se destacan los objetivos futuros, 
entre los que se incluyen la expansión a nuevos dispositivos, la mejora continua del rendimiento y la personalización de 
la interfaz para diferentes clientes.

  \vspace*{25pt}
  \begin{segundoresumo}
    Este traballo presenta o desenvolvemento e implementación dunha plataforma OTT (Over-The-Top) multiplataforma, deseñada 
    para funcionar nunha ampla gama de dispositivos, dende televisores intelixentes ata ordenadores e dispositivos móbiles. 
    A aplicación utiliza unha arquitectura baseada en microservizos, o que garante escalabilidade, flexibilidade e un mantemento 
    sinxelo. O obxectivo principal do proxecto foi crear unha plataforma robusta e adaptable, capaz de integrarse con diversos 
    sistemas operativos e satisfacer as necesidades de múltiples clientes.

    Ao longo do proxecto, abordáronse importantes desafíos técnicos, como a adaptación da aplicación a diferentes capacidades de 
    procesamento e a optimización do rendemento en distintos contornos, tales como WebOS, Tizen e Android TV. Para iso, realizáronse 
    probas exhaustivas, tanto a nivel unitario como de integración e sistema, que permitiron validar a funcionalidade e o rendemento 
    da aplicación en condicións reais de uso.

    Ademais da plataforma OTT, o proxecto inclúe unha aplicación de análise de datos, desenvolvida coa API de Matomo, que permite 
    realizar un seguimento detallado do comportamento dos usuarios. Esta ferramenta facilita a toma de decisións baseada en datos, 
    permitindo aos clientes obter informes personalizados sobre o rendemento dos seus contidos.

    Os resultados obtidos durante o desenvolvemento deste proxecto demostran a capacidade da plataforma para adaptarse ás demandas 
    do mercado e proporcionar unha experiencia de usuario óptima. Finalmente, destácanse os obxectivos futuros, entre os que se 
    inclúen a expansión a novos dispositivos, a mellora continua do rendemento e a personalización da interface para diferentes clientes.
  \end{segundoresumo}
\vspace*{25pt}
\begin{multicols}{2}
\begin{description}
\item [\palabraschaveprincipal:] \mbox{} \\[-20pt]
  \blindlist{itemize}[7] % substitúe este comando por un itemize
                         % que relacione as palabras chave
                         % que mellor identifiquen o teu TFG
                         % no idioma principal da memoria (tipicamente: galego)
\end{description}
\begin{description}
\item [\palabraschavesecundaria:] \mbox{} \\[-20pt]
  \blindlist{itemize}[7] % substitúe este comando por un itemize
                         % que relacione as palabras chave
                         % que mellor identifiquen o teu TFG
                         % no idioma secundario da memoria (tipicamente: inglés)
\end{description}
\end{multicols}

\end{abstract}
\pagestyle{fancy}

%%%%%%%%%%%%%%%%%%%%%%%%%%%%%%%%%%%%%%%%%%%%%%%%%%%%%%%%%%%%%%%%%%%%%%%%%%%%%%%%


 \pagenumbering{roman}
 \setcounter{page}{1}
 \bstctlcite{IEEEexample:BSTcontrol}

 \tableofcontents
 \listoffigures
 \listoftables
 \clearpage
 
 \pagenumbering{arabic}
 \setcounter{page}{1}

 %%%%%%%%%%%%%%%%%%%%%%%%%%%%%%%%%%%%%%%%
 % Capítulos                            %
 %%%%%%%%%%%%%%%%%%%%%%%%%%%%%%%%%%%%%%%%

 \chapter{Introdución}
\label{chap:introducion}

\section{Contexto y motivación}
\label{sec:contexto}
¿Netflix?¿HBO?¿DAZN?¿Amazon Prime Video?¿Disney+? Me atrevería a afirmar que
(casi)todos los lectores de este documento tienen o han tenido acceso a alguna de 
estas aplicaciones en los últimos meses. Diría también que son unas de las principales
aplicaciones en las que uno piensa a la hora de descargar alguna aplicacion en alguno de sus dispositivos.
En muchos de estos dispositivos hasta vienen preinstaladas. Según un estudio
realizado por la Comisión Nacional de los Mercados y la Competencia (CNMC) el 58\% de los hogares 
españoles con acceso a Internet usaba algún servicio de vídeo en streaming a mediados del año 2023, 
frente al 37\% de mediados de 2019. \cite{CNMC}. Otras fuentes situan este porcentaje en un 81\% a principios
de 2023 \cite{Streaming2023} e incluso en un 95\% en junio de 2024 \cite{Streaming2024}. Estas cifras varian un poco dependiendo
de la fuente, pero todas coinciden en algo: el consumo de contenido audiovisual a través de plataformas de streaming
está en auge y ya supera a las plataformas televisivas tradicionales. Leyendo estos estudios y viendo la evolución de
este mercados se entiende el por qué de la constante aparición de nuevas plataformas de este estilo. Desde peliculas y series
hasta deportes, pasando por documentales, cursos, conciertos, etc. Estas plataformas se adaptan a cualquier estilo
de contenido y a cualquier tipo de usuario. 

No es únicamente el número de usuarios el que ha aumentado, sino también el tiempo que pasamos viendo contenido 
en estas plataformas, el número de plataformas distintas a las que accedemos y el número de dispositivos en los
que lo hacemos. La causa de esto es la facilidad de acceso a estas plataformas desde cualquier lugar y dispositivo. 
Y es que, ¿quién no ha empezado a ver una serie en la televisión del salón, ha continuado en la tablet de la cocina y 
ha terminado en el móvil de la cama? ¿O quién no ha parado la serie para cambiar de plataforma y ver un partido 
de fútbol? El consumo de contenido audiovisual nunca había sido tan sencillo y accesible. 

Sencillo para el consumidor claro, pero para llegar a esta versatilidad y facilidad de uso hay un gran trabajo detrás 
para lograr que utilizar estas plataformas sea igual sin importar el tamaño, la marca o el sistema operativo del dispositivo
ni el lugar en el que se encuentre el usuario. 


\section{Objetivo del trabajo}
\label{sec:objetivo}
\subsection{Desarrollo de una aplicación OTT multiplataforma y multicliente}
\label{sec:PlataformaOTT_introduccion}

Aquí surge el objetivo de este trabajo: desarrollar una aplicación multiplataforma para la visualización de contenido
audiovisual en streaming.

En este trabajo se va a abordar todas las fases de desarrollo de una aplicación OTT (Over The Top), con dos características
principales: multiplataforma y multicliente. Multiplataforma porque la aplicación esta orientada a poder adaptarse a cualquier 
dispositivo y sistema operativo. Multicliente porque el código no está pensado únicamente para un cliente en concreto,
sino que está pensado para poder ser utilizado por cualquier cliente que quiera tener su propia plataforma donde mostrar
su contenido.

Estos enfoques suponen varios retos a nivel de desarrollo, los principales: adaptación a las distintas necesidades de cada dispositivo 
donde se quiera dar soporte a la aplicación, con las características y limitaciones que conlleva utilizar las distintas plataformas, tecnologías
y sistemas operativos; y la adaptación a las distintas necesidades de cada cliente, tratando de mantener un equilibrios entre la personalización
de la aplicación y la reutilización del código. 


\subsection{Desarrollo e integración de una aplicación para la recogida y visualización de datos de uso}
\label{sec:Analitica_introduccion}

Cuando se desarrollan aplicaciones comerciales siempre hay un objetivo claro: atraer a los usuarios y mantenerlos el mayor tiempo posible
en la aplicación. Para ello, es necesario conocer a los usuarios, saber qué les gusta, qué no les gusta, qué les interesa, etc. Por ello,
las empresas cada vez más están invirtiendo en herramientas de analítica para poder recoger y analizar los datos de uso de sus aplicaciones.

Complementariamente a la aplicación OTT, en este trabajo se va a abordar el desarrollo de una aplicación de analítica de datos de uso de la
aplicación OTT. Esta aplicación recogerá datos que nos permitirá conocer el comportamiento de los usuarios en la aplicación, para poder
mejorar la experiencia de usuario y adaptar la aplicación a las necesidades de los distintos usuarios de las distintas plataformas.

Este desarrollo incluye el análisis y seleccion de las métricas necesarias para conocer el comportamiento de los usuarios en este tipo de 
plataformas, la integración de la API de analítica en la OTT y el desarrollo de una aplicación de visualización de los datos recogidos.

 \chapter{Fundamentos}
\label{chap:fundamentos}

\section{Introducción}
\label{sec:introduccion_fundamentos}
El desarrollo de aplicaciones multiplataforma representa un gran desafío técnico en un entorno
tan complejo como el tecnológico actualmente, con una variedad tan amplia de dispositivos y plataformas.
Para asegurar el éxito de un proyecto de estas características, es fundamental comprender tanto los principios 
teóricos que permiten crear aplicaciones flexibles y escalables como los fundamentos tecnológicos empleados 
en su desarrollo.

En este capítulo se presentan los fundamentos teóricos y tecnológicos necesarios en el desarrollo de 
las aplicaciones propuestas en este trabajo. Se explorarán conceptos teóricos clave como el funcionamiento 
de aplicaciones OTT, la experiencia de usuario, el diseño de interfaces, la adaptabilidad y la escalabilidad; 
así como las tecnologías esenciales para desarrollar una aplicación que se ajuste tanto a las necesidades del 
cliente como a la plataforma, e integrar una plataforma analítica complementaria.

\section{Fundamentos teóricos}
\label{sec:fundamentos_teoricos}

\subsection{OTT (Over-the-Top)}
\label{sec:fundamentos_teoricos_ott}

\subsubsection{Definición, funcionamiento y clasificación de los servicios OTT}
\label{sec:ott_definicion_funcionamiento_clasificacion}

En radiodifusión, un servicio OTT (siglas en inglés de \textit{over-the-top}) consiste en la transmisión de audio, vídeo y otros 
contenidos a través de internet sin la implicación de los operadores tradicionales en el control o la distribución de los mismos \cite{OTT}.

Estos servicios han transformado radicalmente la forma en que los usuarios consumen contenido audiovisual y cómo los proveedores 
lo distribuyen. La clave de esta transformación radica en la capacidad de acceder a contenidos de manera inmediata, sin necesidad 
de esperar a que se emitan en televisión, radio, cine u otros medios tradicionales. Además, los servicios OTT permiten a los usuarios 
disfrutar de contenidos en cualquier lugar y momento, siempre que cuenten con una conexión a internet.

Su funcionamiento se basa en un catálogo de contenidos que los usuarios pueden explorar según sus preferencias. Cada catálogo 
incluye una variedad de películas, series, documentales y otros tipos de contenido, presentados con descripciones detalladas, 
géneros, calificaciones de usuarios y recomendaciones basadas en el historial de visualización. Los usuarios pueden navegar por 
el catálogo, seleccionar el contenido que desean ver y comenzar la transmisión de manera inmediata, beneficiándose de la capacidad 
de pausar, rebobinar o adelantar el contenido según su conveniencia. Este sistema flexible no solo permite a los usuarios tener control 
total sobre lo que ven y cuándo lo ven, sino que también facilita la entrega personalizada de contenido, adaptándose continuamente a los 
gustos individuales de cada usuario.

El contenido ofrecido por los servicios OTT es muy variado, abarcando desde películas y series hasta música, deportes, noticias, cursos 
educativos y recursos generados por los usuarios. Esta diversidad permite que los servicios OTT se adapten a diferentes preferencias y 
necesidades, ofreciendo una experiencia personalizada para cada tipo de usuario. Además, el acceso a datos masivos sobre los hábitos de 
consumo de los usuarios ha permitido a las plataformas OTT optimizar sus catálogos y estrategias de marketing, logrando una conexión más 
profunda y relevante con su audiencia.

En función del modelo de negocio y la forma en que se ofrece el contenido, los principales tipos de servicios OTT son:
\begin{itemize}
\item \textbf{SVOD (Subscription Video on Demand):} Servicios que ofrecen acceso ilimitado a un catálogo de contenido a cambio de una suscripción mensual o anual.
\item \textbf{AVOD (Advertising Video on Demand):} Servicios que proporcionan contenido gratuito a los usuarios a cambio de la visualización de anuncios publicitarios.
\item \textbf{TVOD (Transactional Video on Demand):} Servicios que permiten a los usuarios alquilar o comprar contenido de manera individual.
\end{itemize}

\subsubsection{Evolución y tendencias de los servicios OTT}
\label{sec:ott_evolucion_tendencias}

El crecimiento de los servicios OTT ha sido impulsado por varios factores clave, como la mejora de las conexiones a internet, que ha permitido una 
transmisión más rápida y de mayor calidad, y los avances tecnológicos en la compresión de video y la optimización de redes, que han mejorado 
significativamente la experiencia del usuario. Estas aplicaciones deben ser diseñadas para escalar eficientemente, manejando grandes volúmenes 
de tráfico durante eventos en vivo o lanzamientos populares, lo que plantea desafíos en términos de rendimiento. Además, la personalización se 
ha convertido en un diferenciador esencial, impulsada por algoritmos de inteligencia artificial que analizan los patrones de consumo para ofrecer 
contenido relevante, mejorando la experiencia del usuario y fomentando la lealtad a la plataforma. La proliferación de dispositivos móviles y la 
preferencia por el contenido instantáneo han consolidado a los servicios OTT como la opción preferida frente a los medios tradicionales, reflejando 
una evolución en los hábitos de consumo hacia la conveniencia y la personalización.

El entorno competitivo también ha impulsado innovaciones en la oferta de servicios OTT. Las empresas han tenido que diferenciarse mediante la 
creación de contenido original exclusivo, el uso de inteligencia artificial para mejorar la personalización y la implementación de modelos de 
negocio híbridos que combinan suscripciones, publicidad y transacciones para maximizar el alcance y la rentabilidad. Estas estrategias no solo 
buscan atraer y retener a los usuarios, sino también garantizar la sostenibilidad del servicio en un mercado cada vez más fragmentado.

\subsection{Arquitectura}
\label{subsec:fundamentos_teoricos_arquitectura}

La arquitectura de software se refiere a los modelos y estándares que sirven de base para el diseño e 
implementación de sistemas de software. Define la estructura, el funcionamiento y la interacción de los 
componentes de software \cite{ArqSoftware}.

El diseño de la arquitectura de un sistema de software para una aplicación OTT es crucial para garantizar 
su correcto funcionamiento, así como una entrega eficiente, escalable y confiable de los contenidos. 
Existen varias opciones para diseñar la arquitectura de una aplicación de estas características, como 
la Arquitectura Monolítica o la Arquitectura Orientada a Servicios (SOA). En este caso, se ha optado por 
una arquitectura basada en microservicios.

\subsubsection{Arquitectura basada en microservicios}
\label{subsec:fundamentos_teoricos_arquitectura_microservicios}

La arquitectura basada en microservicios es un enfoque de diseño de aplicaciones que consiste en un conjunto 
de pequeños servicios, los cuales se ejecutan de forma independiente y se comunican mediante mecanismos 
ligeros (como una API de recursos HTTP, como en este proyecto) \cite{Microservices}. Un microservicio se 
encarga de una única funcionalidad de la aplicación, lo que le permite operar de manera autónoma y comunicarse 
con otros microservicios según sea necesario.

En esta arquitectura, las funcionalidades de la plataforma OTT se distribuyen en diferentes microservicios, 
lo que permite que cada uno sea desarrollado, desplegado y escalado de forma independiente. Esto facilita 
la evolución de la plataforma, ya que cada microservicio puede ser mejorado sin afectar al resto del sistema. 
Además, una buena arquitectura de microservicios mejora la tolerancia a fallos; si uno falla, el resto de la 
plataforma debería continuar funcionando correctamente o verse mínimamente afectado.

Ejemplos de microservicios comunes en una plataforma OTT son:
\begin{itemize}
\item \textbf{Gestión de usuarios:} Encargado de la gestión de usuarios, incluyendo registro, autenticación y autorización.
\item \textbf{Gestión de contenidos:} Responsable de la gestión y almacenamiento de los contenidos, así como de los metadatos y archivos multimedia.
\item \textbf{Orquestador:} Coordina las peticiones de los usuarios a los diferentes microservicios de la plataforma.
\item \textbf{Recomendaciones:} Genera recomendaciones personalizadas para los usuarios.
\end{itemize}

\subsubsection{Componentes de la arquitectura}
\label{subsec:fundamentos_teoricos_arquitectura_componentes}

Además de los microservicios, una plataforma OTT incluye varios componentes clave que conforman la infraestructura 
técnica esencial para su funcionamiento. Estos componentes soportan las operaciones críticas, desde la ingesta y 
distribución de contenido hasta la experiencia de usuario final. Algunos de los componentes más importantes son:

\begin{itemize}
\item \textbf{CDN (Content Delivery Network):} Red de distribución de contenidos que permite la entrega rápida y eficiente de contenidos multimedia a los usuarios finales. La CDN almacena copias de los contenidos en servidores distribuidos geográficamente, reduciendo la latencia y mejorando la velocidad de carga.
\item \textbf{Ingesta y gestión de contenidos:} Componente que recibe, procesa y gestiona todos los contenidos que estarán disponibles en la plataforma. Permite a los administradores subir, editar, etiquetar y organizar los contenidos y sus metadatos para su distribución.
\item \textbf{Gestión de usuarios:} Responsable de todas las funciones relacionadas con los usuarios, incluyendo registro, autenticación, autorización, gestión de perfiles y preferencias.
\item \textbf{Interfaz de usuario:} Es la parte visible de la plataforma, donde los usuarios interactúan con ella. Recibe la información de los microservicios y la presenta de manera amigable.
\item \textbf{Monitorización y análisis:} Componente encargado de monitorizar el rendimiento de la plataforma y de analizar los datos generados por los usuarios, con el fin de mejorar la experiencia de usuario y la eficiencia del sistema.
\item \textbf{Otros componentes:} Incluye componentes de seguridad, monetización, publicidad, gestión de pagos, entre otros, que pueden ser necesarios según las características de la plataforma.
\end{itemize}

\subsection{Experiencia de usuario (UX)}
\label{sec:fundamentos_teoricos_ux}

A diferencia de los medios tradicionales, donde el contenido se emite en un horario fijo y de una forma
especifica, sin darle al usuario la posibilidad de interactuar con el contenido, los servicios OTT permiten
a los usuarios navegar por el catálogo de contenidos, seleccionar lo que desean ver, su información, 
su genetica(temporadas, episodios...) y cualquier otra información relevante, y ver el contenido como y cuando
quieran, sin restricciones de tiempo o lugar. Esta flexibilidad y control que desean los usuarios sobre el contenido
es tanto una de las mejores características de los servicios OTT como uno de los mayores desafíos para los diseñadores.

La experiencia de usuario (UX) \cite{UX} se define como los factores y elementos relativos a la interacción del usuario
con la interfaz de un sistema, dispositivo o aplicación. Estos factores son claves para definir si un usuario disfruta
o no de la experiencia de uso de un producto o servicio. Es fundamental prestar atención a estos factores en el diseño
de cualquier aplicación, y más aún en el caso de una plataforma OTT, donde los usuarios buscan una experiencia fluida,
personalizada, intuitiva y atractiva.

\subsubsection{Principios de diseño de UX}
\label{subsec:fundamentos_teoricos_ux_principios}

El diseño de un UX efectivo se basa en una de principios y prácticas que deben adaptar a cada proyecto, a sus necesidades
y características. En el caso de este tipo de aplicaciones dos de estos prinpios son la usabilidad y la simplicidad.
La plataforma para el consumo de los contenidos debe ser fácil de usar, intuitiva y accesible para todo tipo de usuarios,
independientemente de su nivel de experiencia o conocimientos técnicos. Esto nos asegurará que cualquier usuario podrá 
hacer uso de nuestras plataformas. Otros principios importantes son la consistencia y la adaptabilidad entre dispositivos
y plataformas; la accesibilidad, para garantizar que todos los usuarios puedan disfrutar de la plataforma; y la personalización,
para ofrecer una experiencia única y relevante a cada usuario.

\subsubsection{Elementos clave de la UX en una plataforma OTT}
\label{subsec:fundamentos_teoricos_ux_elementos}

Las plataformas OTT son aplicaciones con una serie de componentes clave que no pueden faltar ya que son necesarios para 
el funcionamiento de la misma. De la misma manera, estos componenetes deben comportarse de una manera determinada. Los
componenetes más importantes son: 

\begin{itemize}
    \item \textbf{Contenidos} Elemento central de la plataforma sobre el que gira toda la experiencia de usuario. Los contenidos
    deben ser fáciles de encontrar, navegar y consumir, y deben estar presentados de forma atractiva y organizada.
    \item \textbf{Agrupaciones de contenidos} Ya sea en listas, series, temporadas, categorías, géneros, etc. Las agrupaciones
    de contenidos permiten a los usuarios explorar y descubrir nuevos contenidos de forma sencilla y atractiva.
    \item \textbf{Metadata} información sobre cualquiera de los elementos que conforman la plataforma: descripciones, fechas,
    títulos, actores, etc. Aporta al usuario el contexto necesario para entender y valorar los contenidos.
    \item \textbf{Páginas} Cada menú, serie, temporada, episodio, etc. debe tener su propia página con la información y funcionalidades
    necesarias para hacer uso de la plataforma.
    \item \textbf{Botones y controles} Los botones y controles de reproducción, pausa, adelanto, retroceso, etc. deben ser intuitivos
    y accesibles para que los usuarios puedan interactuar con el contenido de forma sencilla.
\end{itemize}


\input{contido/Fundamentos/Teoricos/Esc_Adapt.text}
 \chapter{Contido demostrativo}
\label{chap:demo}

\lettrine{E}{ntre} a introdución e as conclusións, o documento conterá
tantos capítulos como sexa preciso, sempre con coidado de non rebasar
o límite de 80 páxinas fixado polo regulamento de TFGs.

Empregaremos éste de xeito demostrativo, para ilustrar o uso de
elementos habituais que poidan ser de utilidade\footnote{Por exemplo,
  isto é unha nota a pé de páxina.}.

\section{Inclusión de imaxes}

Se precisamos imaxes no noso documento, incluirémolas do xeito que se
indica na figura~\ref{fig:exemplo} (páxina~\pageref{fig:exemplo}). Se
o facemos así, \LaTeX ubicará cada imaxe no mellor lugar posible,
lugar que pode variar a medida que o documento vaia crecendo coa
inclusión de máis texto e outros elementos (máis imaxes, táboas,
etc.).

\begin{figure}[hp!]
  \centering
  \includegraphics[width=0.75\textwidth]{imaxes/udc.png}
  \caption{Pé de imaxe descritivo}
  \label{fig:exemplo}
\end{figure}

Recoméndase almacenar os ficheiros gráficos no directorio
\texttt{imaxes}.

\subsection{Inclusión de varias sub-imaxes}

Se precisamos inserir imaxes relacionadas, pode ser apropiado
incluílas como sub-figuras, do xeito que se pode apreciar na
figura~\ref{fig:exemplo-subfiguras} (páxina~\pageref{fig:exemplo-subfiguras})
coas imaxes~\ref{fig:subfigura-rotada} e~\ref{fig:subfigura-deformada}.
Como se pode ver nos exemplos desta sección, sempre é recomendable
referirse ás imaxes (ou táboas e outros elementos \emph{flotantes},
que se demostrarán nas seccións seguintes deste capítulo demostrativo)
pola súa referencia, xa que dese xeito non dependemos de onde
queden ubicados os elementos en cuestión.

\begin{figure}[hp!]
  \centering
  \begin{subfigure}[c]{0.3\textwidth}
    \includegraphics[angle=45,width=\textwidth]{imaxes/udc.png}
    \caption{Pé de subimaxe rotada}
    \label{fig:subfigura-rotada}
  \end{subfigure}
  \hspace{0.1\textwidth}
  \begin{subfigure}[c]{0.3\textwidth}
    \includegraphics[width=\textwidth,height=3cm]{imaxes/udc.png}
    \caption{Pé de subimaxe deformada}
    \label{fig:subfigura-deformada}
  \end{subfigure}
  \caption{Pé de imaxe xeral}
  \label{fig:exemplo-subfiguras}
\end{figure}

\section{Inclusión de táboas}

Se precisamos táboas no noso documento, incluirémolas do xeito que se
indica na táboa~\ref{tab:exemplo} (páxina~\pageref{tab:exemplo}). Se
o facemos así, \LaTeX ubicará cada táboa no mellor lugar posible,
lugar que pode variar a medida que o documento vaia crecendo coa
inclusión de máis texto e outros elementos (máis imaxes, táboas,
etc.).

\begin{table}[hp!]
  \centering
  \rowcolors{2}{white}{udcgray!25}
  \begin{tabular}{c|c}
  \rowcolor{udcpink!25}
  \textbf{Título de columna} & \textbf{Outro título de columna} \\\hline
  \textit{Título de fila} & Contido da cela \\
  \textit{Título de fila} & Contido da cela \\
  \textit{Título de fila} & Contido da cela \\
  \textit{Título de fila} & Contido da cela \\
  \textit{Título de fila} & Contido da cela \\
  \textit{Título de fila} & Contido da cela \\
  \end{tabular}
  \caption{Pé de táboa descritivo}
  \label{tab:exemplo}
\end{table}

\subsection{Inclusión de táboas longas}

Para táboas longas que ocupan varias páxinas, como é o caso da \ref{tab:longa}
(páxina~\pageref{tab:longa}), recoméndase o uso do paquete \texttt{lontable},
incluído xa entre os paquetes recomendados no ficheiro raíz do proxecto
(\verb+memoria_tfg.tex+).

{\rowcolors{2}{white}{udcgray!25}
\begin{longtable}{l|r|c}
  \caption{Pé descritivo dunha táboa longa}
  \label{tab:longa} \\

  \rowcolor{udcpink!25}
  \textbf{Primeira columna} & \textbf{Segunda columna} & \textbf{Terceira columna} \\\hline
  \endfirsthead

  \multicolumn{3}{c}{\tablename\ \thetable{} -- {\small \textit{(vén da páxina anterior)}}} \\
  \rowcolor{udcpink!25}
  \textbf{Primeira columna} & \textbf{Segunda columna} & \textbf{Terceira columna} \\\hline
  \endhead

  \multicolumn{3}{c}{\dotfill{\small \textit{(continúa na páxina seguinte)}}\dotfill} \\
  \endfoot

  \endlastfoot

  Texto de exemplo & abcdef ghjijklmn & 123.456778 \\
  Texto de exemplo & abcdef ghjijklmn & 123.456778 \\
  Texto de exemplo & abcdef ghjijklmn & 123.456778 \\
  Texto de exemplo & abcdef ghjijklmn & 123.456778 \\
  Texto de exemplo & abcdef ghjijklmn & 123.456778 \\
  Texto de exemplo & abcdef ghjijklmn & 123.456778 \\
  Texto de exemplo & abcdef ghjijklmn & 123.456778 \\
  Texto de exemplo & abcdef ghjijklmn & 123.456778 \\
  Texto de exemplo & abcdef ghjijklmn & 123.456778 \\
  Texto de exemplo & abcdef ghjijklmn & 123.456778 \\
  Texto de exemplo & abcdef ghjijklmn & 123.456778 \\
  Texto de exemplo & abcdef ghjijklmn & 123.456778 \\
  Texto de exemplo & abcdef ghjijklmn & 123.456778 \\
  Texto de exemplo & abcdef ghjijklmn & 123.456778 \\
  Texto de exemplo & abcdef ghjijklmn & 123.456778 \\
  Texto de exemplo & abcdef ghjijklmn & 123.456778 \\
  Texto de exemplo & abcdef ghjijklmn & 123.456778 \\
  Texto de exemplo & abcdef ghjijklmn & 123.456778 \\
  Texto de exemplo & abcdef ghjijklmn & 123.456778 \\
  Texto de exemplo & abcdef ghjijklmn & 123.456778 \\
  Texto de exemplo & abcdef ghjijklmn & 123.456778 \\
  Texto de exemplo & abcdef ghjijklmn & 123.456778 \\
  Texto de exemplo & abcdef ghjijklmn & 123.456778 \\
  Texto de exemplo & abcdef ghjijklmn & 123.456778 \\
  Texto de exemplo & abcdef ghjijklmn & 123.456778 \\
  Texto de exemplo & abcdef ghjijklmn & 123.456778 \\
  Texto de exemplo & abcdef ghjijklmn & 123.456778 \\
  Texto de exemplo & abcdef ghjijklmn & 123.456778 \\
  Texto de exemplo & abcdef ghjijklmn & 123.456778 \\
  Texto de exemplo & abcdef ghjijklmn & 123.456778 \\
  Texto de exemplo & abcdef ghjijklmn & 123.456778 \\
  Texto de exemplo & abcdef ghjijklmn & 123.456778 \\
  Texto de exemplo & abcdef ghjijklmn & 123.456778 \\
  Texto de exemplo & abcdef ghjijklmn & 123.456778 \\
  Texto de exemplo & abcdef ghjijklmn & 123.456778 \\
  Texto de exemplo & abcdef ghjijklmn & 123.456778 \\
  Texto de exemplo & abcdef ghjijklmn & 123.456778 \\

\end{longtable}
} % alcance do rowcolors que precede á táboa longa


\subsection{Inclusión de táboas con celas que ocupan varias columnas ou filas}

En ocasións pode resultar de interese incluír nunha táboa unha cela que se estenda
a través de varias columnas, como ocorre na táboa~\ref{tab:exemplocolumnas}
(páxina~\pageref{tab:exemplocolumnas}).

\begin{table}[hp!]
  \centering
  \rowcolors{2}{white}{udcgray!25}
  \begin{tabular}{c|c|c}
  \rowcolor{udcpink!25}
  \multicolumn{3}{c}{\textbf{Cela en varias columnas}} \\\hline
  \rowcolor{udcpink!25}
  \textbf{Título de columna} & \textbf{Outro título de columna} & \textbf{Outro título máis} \\\hline
  \textit{Título de fila}    & Contido da cela                  & Contido da cela \\
  \textit{Título de fila}    & Contido da cela                  & Contido da cela \\
  \textit{Título de fila}    & \multicolumn{2}{c}{Contido da cela múltiple} \\
  \textit{Título de fila}    & Contido da cela                  & Contido da cela \\
  \end{tabular}
  \caption{Pé de táboa descritivo (táboa con celas que ocupan varias columnas)}
  \label{tab:exemplocolumnas}
\end{table}

Tamén pode resultar necesario facer o propio mais en varias filas da mesma columna,
como ocorre na táboa~\ref{tab:exemplofilas} (páxina~\pageref{tab:exemplofilas}).
Para isto é preciso o paquete \texttt{multirow}, incluído entre os recomendados no
ficheiro raíz do proxecto (\verb+memoria_tfg.tex+).

O uso de celas multifila requerirá do xuste da coloración das filas, a fin de manter
a coherencia entre o contido e o continente. Así, no canto de usar un único comando
\verb+rowcolors+ para indicar a alternancia en toda a táboa, usaremos o comando
\verb+rowcolor+ antes dunha fila que queiramos colorear, e o comando \verb+cellcolor+
dentro dunha cela que queiramos colorear.

\begin{table}[hp!]
  \centering
  \begin{tabular}{c|c}
  \rowcolor{udcpink!25}
  \textbf{Título de columna} & \textbf{Outro título de columna} \\\hline
  \multirow{2}{*}{\textit{Título de fila}} & \cellcolor{udcgray!25} Contido da cela \\
                                           & Contido da cela \\
  \rowcolor{udcgray!25}
  \textit{Título de fila}                  & Contido da cela \\
  \multirow{3}{*}{\textit{Título de fila}} & Contido da cela \\
                                           & \cellcolor{udcgray!25} Contido da cela \\
                                           & Contido da cela \\
  \rowcolor{udcgray!25}
  \textit{Título de fila}                  & Contido da cela \\
  \end{tabular}
  \caption{Pé de táboa descritivo (táboa con celas que ocupan varias filas)}
  \label{tab:exemplofilas}
\end{table}

Por suposto, pódense combinar nunha mesma táboa os dous tipos de celas (as que se 
estenden máis dunha fila e máis dunha columna), como na táboa~\ref{tab:exemplofilasecolumnas}
(páxina~\pageref{tab:exemplofilasecolumnas}).

\begin{table}[hp!]
  \centering
  \begin{tabular}{c|c|c}
  \rowcolor{udcpink!25}
  \multicolumn{3}{c}{\textbf{Cela en varias columnas}} \\\hline
  \rowcolor{udcpink!25}
  \textbf{Título de columna}               & \textbf{Outro título de columna}             & \textbf{Outro título máis} \\\hline
  \multirow{2}{*}{\textit{Título de fila}} & \cellcolor{udcgray!25} Contido da cela       & \cellcolor{udcgray!25} Contido da cela \\
                                           & Contido da cela                              & Contido da cela \\
  \rowcolor{udcgray!25}
  \textit{Título de fila}                  & \multicolumn{2}{c}{Contido da cela múltiple} \\
  \multirow{3}{*}{\textit{Título de fila}} & Contido da cela                              & Contido da cela \\
                                           & \multicolumn{2}{c}{\cellcolor{udcgray!25} Contido da cela múltiple} \\
                                           & \multicolumn{2}{c}{Contido da cela múltiple} \\
  \rowcolor{udcgray!25}
  \textit{Título de fila}                  & Contido da cela                              & Contido da cela \\
  \end{tabular}
  \caption{Pé de táboa descritivo (táboa con celas que ocupan varias columnas)}
  \label{tab:exemplofilasecolumnas}
\end{table}

\section{Inclusión de código fonte}

Se precisamos incluír fragmentos de código fonte, podemos facelo, por exemplo, da
seguinte maneira:

\begin{lstlisting}[language=C]
#include <stdio.h>
#define N 10

int main()
{
  int i;

  // Isto é un comentario
  puts("Ola, mundo!");

  for (i = 0; i < N; i++)
  {
    puts("LaTeX é a ferramenta de edición ideal para profesionais da informática!");
  }

  return 0;
}
\end{lstlisting}

\section{Uso da relación de acrónimos e do glosario}

Os acrónimos edítanse no ficheiro \texttt{bibliografia/acronimos.tex}
e úsanse empregando a orde \texttt{acrlong} para obter o termo
completo (deste xeito: \acrlong{erlang}), a orde \texttt{acrshort}
para obter o acrónimo (deste xeito: \acrshort{erlang}). A primeira vez
que usamos un termo con acrónimo no documento é recomendable usar orde
\texttt{acrfull} (que produce ambas versións á vez:
\acrfull{erlang}). Os acrónimos que non se usan no documento, non
aparecen na relación que se xerar na versión PDF.

Pola súa banda, os termos do glosario edítanse no ficheiro
\texttt{bibliografia/glo\-sa\-rio.tex} e úsanse empregando a orde
\texttt{gls} (deste xeito, \gls{bytecode}) ou \texttt{Gls} (deste
xeito, \Gls{bytecode}). Ao igual que os acrónimos, os termos que non
se usan no documento, non aparecen na relación que se xera na versión
PDF.

%\include{contido/...}
 \chapter{Conclusións}
\label{chap:conclusions}

\lettrine{D}{erradeiro} capítulo da memoria, onde se presentará a
situación final do traballo, as leccións aprendidas, a relación coas
competencias da titulación en xeral e a mención en particular,
posibles liñas futuras,\dots

\Blindtext


 %%%%%%%%%%%%%%%%%%%%%%%%%%%%%%%%%%%%%%%%
 % Apéndices, glosarios e bibliografía  %
 %%%%%%%%%%%%%%%%%%%%%%%%%%%%%%%%%%%%%%%%

 \appendix
 \appendixpage
 \chapter{Material adicional}
\label{chap:adicional}

\lettrine{E}{xemplo} de capítulo con formato de apéndice, onde se pode
incluír material adicional que non teña cabida no corpo principal do
documento, suxeito á limitación de 80 páxinas establecida no
regulamento de TFGs.

\Blindtext

%\include{anexos/...}

 \printglossary[type=\acronymtype,title=\nomeglosarioacronimos]
 \printglossary[title=\nomeglosariotermos]

 \bibliographystyle{IEEEtranN}
 \bibliography{\bibconfig,bibliografia/bibliografia}
 \clearpage
 
\end{document}

%%%%%%%%%%%%%%%%%%%%%%%%%%%%%%%%%%%%%%%%%%%%%%%%%%%%%%%%%%%%%%%%%%%%%%%%%%%%%%%%
